四种主模型中，随机森林的汇总 Macro-F1 最高（0.9279），Accuracy 为 0.9236。融合模型相对随机森林的汇总 F1 差值为 -0.56 个百分点。表\ref{tab:main}同时展示个人均值与离散程度，用于区分总体窗口表现与个体迁移波动。

按个人实际出现类别 F1 进行配对，融合减随机森林的均值为 -0.49 个百分点，受试者重采样百分位区间为 [-2.63, 0.99] 个百分点，探索性符号翻转双侧 $p=0.8750$。这些统计只基于八个配对受试者，不能把区间或单个最好分数当作大样本人群中的保证。

融合与早期 CNN 的输入表示和容量均不同；即使融合成绩较高，也不能只据此证明多分支结构本身有效。现有实验没有单独删除融合模型的统计分支或保持等参数量，因而只能比较完整系统。